%% Included in decks: slidemage-manual, slidemage-manual-copy
%%

%% Translation of: 1us-EN-Installing-LaTeX
%%

% SlideMage manual, chapter "Installing SlideMage": the TeX distribution.
% Not a click-by-click guide: where to download, which choices matter,
% what to do afterwards.
\begin{standardframe}{LaTeX installieren}
	\small
	\textbf{Diese Folie gilt nur für Windows.}
	SlideMage übersetzt jedes Deck mit LuaLaTeX über latexmk und braucht daher eine
	TeX-Distribution. MiKTeX ist die kleinere Installation: Es lädt jedes
	Paket erst, wenn ein Deck es zum ersten Mal braucht. TeX Live installiert alles auf einmal.
	\par\medskip
	\begin{columns}[T,onlytextwidth]
		\begin{column}{0.56\linewidth}
			\shrinklists{1}
			\textbf{MiKTeX} (\SMCode{miktex.org/download}, der Basic Installer)
			\begin{itemize}
				\item „Only for me“ braucht keine Administratorrechte
				\item \emph{Install missing packages on-the-fly} auf \emph{Yes} stellen, sonst fragt jeder erste Übersetzungslauf nach
				\item Danach: MiKTeX Console $\rightarrow$ \emph{Updates} $\rightarrow$ \emph{Check for updates}, dann \emph{Update now}
				\item latexmk braucht Perl: Git for Windows installieren (\SMCode{git-scm.com}, Standardoptionen); SlideMage findet dessen Perl selbst
			\end{itemize}
		\end{column}
		\begin{column}{0.4\linewidth}
			\shrinklists{1}
			\textbf{TeX Live} (\SMCode{tug.org/texlive})
			\begin{itemize}
				\item \SMCode{install-tl-windows.exe}, Schema \emph{full} (etwa 8\,GB, Perl inklusive)
			\end{itemize}
			\medskip
			\textbf{In beiden Fällen}
			\begin{itemize}
				\item Der erste Übersetzungslauf lädt Schriften und Pakete: Er braucht Internet und dauert einige Minuten
				\item Im Terminal testen: \SMCode{lualatex --version}, \SMCode{latexmk -v}
			\end{itemize}
		\end{column}
	\end{columns}
\end{standardframe}
